\section{生物学示例族 I：睡眠--清醒切换}

对本手稿而言，可用的最强生物学转变模型，仍然是第~\ref{ch:consciousness-biology} 章中的睡眠转变来源。
它在这里重要，是因为它为受约束地使用分岔与阈值语言提供了正当性。

\begin{discussion}
睡眠--清醒切换在教学上有三点价值：
\begin{enumerate}[nosep]
  \item 它是真正的生物学机制转变，而不是某种自助式轶事；
  \item 它表明状态变化可以是突变的、有模式的，并且对参数敏感；
  \item 它表明关键数学问题往往不是“当前有多少努力”，而是“哪一个稳定性条件正在被改变”。
\end{enumerate}
\end{discussion}

这\emph{并不}意味着阅读一条定理与入睡是同一个机制。
证据车道仍然保持清晰。
强车道支持在生物学中普遍正当地使用“机制切换”语言；
迁移到行为情境中，仍然属于\textbf{中等}强度的解释性推进。
但这一推进仍然有价值，因为它鼓励读者用阈值、时机与盆地几何来思考，而不只是诉诸欲望强弱。